\section*{Capítulo \ref{ch:g11:redox} --- Oxidación y reducción}

\begin{solution}{exo:g11:redox:1}
\omterm{def:g11:redox:oxidant}{Oxidante}: una especie capaz de ganar \omterm{def:g9:inside-the-atom:nucleus}{electrones}. \omterm{def:g11:redox:oxidant}{Reductor}: una especie
capaz de ceder \omterm{def:g9:inside-the-atom:nucleus}{electrones}. \omterm{def:g11:redox:oxidation}{Oxidación}: una pérdida de \omterm{def:g9:inside-the-atom:nucleus}{electrones}.
\omterm{def:g11:redox:oxidation}{Reducción}: una ganancia de \omterm{def:g9:inside-the-atom:nucleus}{electrones}.
\end{solution}

\begin{solution}{exo:g11:redox:2}
\ce{Fe^{2+}}/\ce{Fe}; \ce{Ag+}/\ce{Ag}; \ce{I2}/\ce{I-};
\ce{Fe^{3+}}/\ce{Fe^{2+}}.
\end{solution}

\begin{solution}{exo:g11:redox:3}
\ce{Ag+ + e- <=> Ag}; \ce{Al^{3+} + 3e- <=> Al};
\ce{Fe^{3+} + e- <=> Fe^{2+}}; \ce{Cl2 + 2e- <=> 2Cl-}.
\end{solution}

\begin{solution}{exo:g11:redox:4}
Una \omterm{def:g11:redox:oxidation}{oxidación} (se pierden \omterm{def:g9:inside-the-atom:nucleus}{electrones}); el hierro actúa como \omterm{def:g11:redox:oxidant}{reductor}.
\end{solution}

\begin{solution}{exo:g11:redox:5}
El cinc se oxida y es el \omterm{def:g11:redox:oxidant}{reductor}; el \ce{Cu^{2+}} se reduce y es el
\omterm{def:g11:redox:oxidant}{oxidante}. Dos \omterm{def:g9:inside-the-atom:nucleus}{electrones} por \omterm{def:g7:atoms-and-molecules:atom}{átomo} de cinc.
\end{solution}

\begin{solution}{exo:g11:redox:6}
Se parte de \ce{H2O2 -> H2O}: % ce-unbalanced-ok (step 1 of the method)
oxígeno: dos O a la izquierda, así que \ce{2H2O} a la derecha;
hidrógeno: 2 a la izquierda, 4 a la derecha, así que \ce{2H+} a la
izquierda; carga: $+2$ a la izquierda, 0 a la derecha, así que dos
\omterm{def:g9:inside-the-atom:nucleus}{electrones} a la izquierda:
\ce{H2O2 + 2H+ + 2e- <=> 2H2O}.
\end{solution}

\begin{solution}{exo:g11:redox:7}
\ce{I2 + 2e- -> 2I-} y \ce{2S2O3^{2-} -> S4O6^{2-} + 2e-}; dos \omterm{def:g9:inside-the-atom:nucleus}{electrones}
cada una, así que \ce{I2 + 2S2O3^{2-} -> 2I- + S4O6^{2-}}.
\end{solution}

\begin{solution}{exo:g11:redox:8}
\ce{Cu -> Cu^{2+} + 2e-} y $2\times$ (\ce{Ag+ + e- -> Ag}):
\ce{Cu + 2Ag+ -> Cu^{2+} + 2Ag}. Los \omterm{def:g9:ions:ion}{iones} \ce{Cu^{2+}} formados son
azules.
\end{solution}

\begin{solution}{exo:g11:redox:9}
\ce{C6H8O6 -> C6H6O6 + 2H+ + 2e-} y \ce{I2 + 2e- -> 2I-}:
\ce{C6H8O6 + I2 -> C6H6O6 + 2I- + 2H+}. La vitamina C es el \omterm{def:g11:redox:oxidant}{reductor}.
\end{solution}

\begin{solution}{exo:g11:redox:10}
\ce{Al -> Al^{3+} + 3e-} da 3 \omterm{def:g9:inside-the-atom:nucleus}{electrones}, y \ce{Cu^{2+} + 2e- -> Cu}
toma 2; el menor número común es 6, así que la primera se multiplica por
2 y la segunda por 3: \ce{2Al + 3Cu^{2+} -> 2Al^{3+} + 3Cu}.
\end{solution}

\begin{solution}{exo:g11:redox:11}
\ce{H2O2 + 2H+ + 2e- -> 2H2O} y $2\times$ (\ce{Fe^{2+} -> Fe^{3+} + e-}):
\ce{H2O2 + 2H+ + 2Fe^{2+} -> 2H2O + 2Fe^{3+}}.
\end{solution}

\begin{solution}{exo:g11:redox:12}
$n(\ce{Zn}) = 1.30/65.4 = \qty{1.99e-2}{mol}$. Se deposita un \omterm{def:g7:atoms-and-molecules:atom}{átomo} de
cobre por cada \omterm{def:g7:atoms-and-molecules:atom}{átomo} de cinc oxidado, así que
$n(\ce{Cu}) = \qty{1.99e-2}{mol}$ y
$m(\ce{Cu}) = 1.99\times 10^{-2} \times 63.5 = \qty{1.26}{g}$.
\end{solution}

\begin{solution}{exo:g11:redox:13}
\ce{Cr2O7^{2-} + 14H+ + 6e- -> 2Cr^{3+} + 7H2O} y
$6\times$ (\ce{Fe^{2+} -> Fe^{3+} + e-}):
\ce{Cr2O7^{2-} + 14H+ + 6Fe^{2+} -> 2Cr^{3+} + 6Fe^{3+} + 7H2O}.
Seis \ce{Fe^{2+}} por \omterm{def:g9:ions:ion}{ion} dicromato: $6 \times 0.010 = \qty{0.060}{mol}$.
\end{solution}

\begin{solution}{exo:g11:redox:14}
\ce{2ClO- + 4H+ + 2e- -> Cl2 + 2H2O} y \ce{2Cl- -> Cl2 + 2e-}; sumando
y dividiendo entre 2: \ce{ClO- + Cl- + 2H+ -> Cl2 + H2O}. Un ácido aporta
los \omterm{def:g9:ions:ion}{iones} \ce{H+}, y la reacción libera dicloro, un gas tóxico: de ahí
la advertencia de no \omterm{def:g2:mixing-and-dissolving:mix}{mezclar} nunca los dos productos.
\end{solution}

\begin{solution}{exo:g11:redox:15}
Los \omterm{def:g9:inside-the-atom:nucleus}{electrones} no pueden moverse libremente por la disolución
(\cref{prop:g11:redox:no-free-electrons}): los que pierden los \omterm{def:g7:atoms-and-molecules:atom}{átomos} de
cinc solo los puede tomar un \omterm{def:g9:ions:ion}{ion} \ce{Cu^{2+}} que toque el \omterm{def:g9:periodic-table-first-look:metal}{metal}. Por
tanto, el \omterm{def:g7:atoms-and-molecules:atom}{átomo} de cobre se forma en la superficie de la lámina, y se
queda allí.
\end{solution}

\begin{solution}{pb:g11:redox:1}
\textbf{1.} \ce{Cr2O7^{2-}}/\ce{Cr^{3+}} y \ce{CH3COOH}/\ce{CH3CH2OH}.

\textbf{2.} El etanol se oxida: es el \omterm{def:g11:redox:oxidant}{reductor}.

\textbf{3.} \ce{Cr2O7^{2-} + 14H+ + 6e- <=> 2Cr^{3+} + 7H2O}.

\textbf{4.} El carbono está ajustado (dos en cada miembro). Oxígeno:
dos a la izquierda, uno a la derecha, así que \ce{H2O} a la derecha.
Hidrógeno: cuatro a la izquierda, $6 + 2 = 8$ a la derecha, así que
\ce{4H+} a la izquierda. Carga: $+4$ a la izquierda, 0 a la derecha,
así que cuatro \omterm{def:g9:inside-the-atom:nucleus}{electrones} a la izquierda:
\ce{CH3COOH + 4H+ + 4e- <=> CH3CH2OH + H2O}.

\textbf{5.} Doce \omterm{def:g9:inside-the-atom:nucleus}{electrones}: $2\times$ la primera y $3\times$ la segunda
invertida:
\ce{2Cr2O7^{2-} + 3CH3CH2OH + 16H+ -> 4Cr^{3+} + 3CH3COOH + 11H2O}.

\textbf{6.} Los \omterm{def:g9:ions:ion}{iones} \ce{Cr2O7^{2-}}, naranjas, se reducen a \omterm{def:g9:ions:ion}{iones}
\ce{Cr^{3+}}, verdes, allí donde llega el etanol.

\textbf{7.} $M = 2\times 39.1 + 2\times 52.0 + 7\times 16.0 =
\qty{294.2}{g/mol}$.

\textbf{8.} $n = 2.0\times 10^{-3}/294.2 = \qty{6.8e-6}{mol}$.

\textbf{9.} Tres \omterm{def:g7:atoms-and-molecules:molecule}{moléculas} de etanol por cada dos \omterm{def:g9:ions:ion}{iones} dicromato:
$n = \tfrac{3}{2}\times 6.8\times 10^{-6} = \qty{1.0e-5}{mol}$.

\textbf{10.} $M(\ce{C2H6O}) = \qty{46.0}{g/mol}$, así que
$m = 1.02\times 10^{-5}\times 46.0 = \qty{4.7e-4}{g}$, unos
\qty{0.47}{mg}.

\textbf{11.} $0.25/0.47 \approx 0.53$: alrededor del \qty{53}{\percent}
del dicromato.

\textbf{12.} $0.53 \times 10 \approx \qty{5.3}{cm}$.

\textbf{13.} $\qty{0.25}{mg} \times 2100 = \qty{525}{mg}$ por litro de
sangre, unos \qty{0.53}{g/L}.

\textbf{14.} El aliento encuentra primero los granos de la entrada; el
dicromato de esa zona está en exceso y retiene todo el etanol, así que
la reacción es completa en un primer tramo del tubo, que se vuelve
totalmente verde, antes de que llegue etanol más lejos.

\textbf{15.} La reacción es total: los \omterm{def:g9:ions:ion}{iones} \ce{Cr^{3+}}, verdes, no
vuelven a ser dicromato, así que un tubo usado no puede volver a medir.

% ledger: ghs:K2Cr2O7 (H350 among its hazard statements)
\textbf{16.} Es de toxicidad aguda (GHS06) y presenta un grave peligro
para la salud (GHS08: entre sus indicaciones de peligro figura ``puede
provocar cáncer''), además de ser corrosivo y tóxico para los organismos
acuáticos; un aparato que usa el público no debe contenerlo.

\textbf{17.} \ce{O2 + 4H+ + 4e- -> 2H2O} y
\ce{CH3CH2OH + H2O -> CH3COOH + 4H+ + 4e-}:
\ce{CH3CH2OH + O2 -> CH3COOH + H2O}.

\textbf{18.} Cuatro \omterm{def:g9:inside-the-atom:nucleus}{electrones}.

\textbf{19.} $n = 0.25\times 10^{-3}/46.0 = \qty{5.4e-6}{mol}$ de
etanol; \omterm{def:g9:inside-the-atom:nucleus}{electrones}: $4n = \qty{2.2e-5}{mol}$, es decir,
$2.2\times 10^{-5}\times 6.02\times 10^{23} = \num{1.3e19}$ \omterm{def:g9:inside-the-atom:nucleus}{electrones}.

\textbf{20.} $Q = 1.3\times 10^{19}\times 1.60\times 10^{-19} \approx
\qty{2.1}{C}$. Cada \omterm{def:g7:atoms-and-molecules:molecule}{molécula} de etanol envía exactamente cuatro
\omterm{def:g9:inside-the-atom:nucleus}{electrones} por el cable, así que la carga es proporcional a la cantidad
de etanol: medirla es medir el etanol.
\end{solution}
